\section{Conclusões}\label{sec:conclusoes}

Apresentamos o GraphCube Lab, uma plataforma que parte de uma visualização existente --- o diagrama de círculos do GraphCube --- e a insere num arcabouço matemático e computacional verificável. Mostramos que o diagrama original é um grafo de incidência entre posições de adesivos e camadas, distinguimos formalmente essa estrutura do grafo de estados e dos grafos de Cayley, e implementamos as três representações de modo que cada elemento visual tenha interpretação precisa. O núcleo matemático, derivado de uma única descrição geométrica e de uma convenção de composição explícita, reproduz resultados conhecidos (ordens de grupos por Schreier--Sims e a distribuição de distâncias do \cube{2}). Os algoritmos de busca são instrumentados e visualizados a partir de seus registros reais de execução, e os experimentos são reprodutíveis a partir de sementes, metadados e um único comando.

Os experimentos quantificaram fenômenos esperados pela teoria com dados próprios: a saturação da distância real em função do comprimento do embaralhamento, a redução do crescimento por camada da busca bidirecional para aproximadamente a raiz quadrada, o ganho de combinar \textit{pattern databases} pelo máximo, o compromisso tempo--memória entre A$^*$ e IDA$^*$, o contraste entre a solução ótima viável apenas para estados próximos no \cube{3} e as soluções válidas do método de duas fases, e o efeito do conjunto gerador sobre diâmetro, cintura e simetrias de grafos de Cayley.

\paragraph{Trabalhos futuros.} (i) PDBs maiores e aditivas~\cite{korf1997,felner2004}, e redução por simetria~\cite{rokicki2014}, para ampliar o alcance das buscas ótimas no \cube{3}; (ii) cubos $n\times n\times n$ maiores, cuja geometria já é genérica no núcleo; (iii) grafos de Schreier (quocientes por subgrupos) completos e visualizáveis, como o das orientações do \cube{2}; (iv) estimativas de mistura de passeios aleatórios em grafos de Cayley; (v) avaliação da interface com estudantes, como ferramenta didática de teoria dos grupos.

\paragraph{Disponibilidade.} Código-fonte (licença MIT), dados brutos, scripts de análise e o código \LaTeX{} deste artigo estão no repositório do projeto, com instruções de reprodução no \texttt{README}.
